\chapter*{Словарь терминов}             % Заголовок
\addcontentsline{toc}{chapter}{Словарь терминов}  % Добавляем его в оглавление

\textbf{верхний фрейм СП} - рама, на которой крепится объект управления (кабина, кресло оператора, пр.)

\textbf{нижний фрейм СП} - рама, на которой размещаются двигатели и центральная опора СП

\textbf{ограничительная характеристика} - поверхность, ограничивающая рабочую зону СП

\textbf{рабочая зона СП} - область пространства, в котором может находиться центр масс подвижной платформы (либо область допустимых значений угловых и поступательных координат платформы)

\textbf{сервопривод} - привод, замкнутый обратной связью по положению выходного звена исполнительного механизма

\textbf{преобразователь частоты / частотный преобразователь} - устройство, предназначенное для управления частотой подводимого напряжения

\textbf{акселерационное воздействие / эффект движения} - поступательное или угловое ускорение, возникающее при движении объекта и воспринимаемое вестибулярным аппаратом человека

\textbf{приводное звено} - входное звено механизма, приводимое в движение с помощью мотор-редуктора.

\textbf{транспортная задержка} - задержка при передаче управляющего воздействия от регулятора к объекту или задержка получения значений управляемых переменных от датчиков.

\textbf{каскадное регулирование} - регулирование, в котором два или больше контуров регулирования соединены так, чтобы выход одного регулятора корректировал уставку другого регулятора.
